\section{Evaluation DAG}

Once the initial evaluation (Fig.~\ref{fig:evaluation}) is complete and its
parametric specification \maxime{init-spec} built, an evaluation directed
acyclic graph (or DAG) must be issued for each registered persistent name before
a parametric specification can be re-evaluated. %
An Evaluation DAG traces the history of each topological parameter back to the
first created orbits it originates from, thus allowing the matching of the
corresponding topological parameter at re-evaluation time. %

\subsection{Evaluation}

\begin{figure}[ht]
  \centering
  \includegraphics[width=0.3\textwidth]{pn3}
  \caption{Evaluation DAG of the Persistent Name $PN_{3}$}
  \label{fig:eval-dag-pn3}
\end{figure}

An evaluation DAG is a data structure used to store the history of a topological
parameter and is built with only the rules and nodes referenced in a persistent
names. %

Since a persistent name only references a topological parameter at the moment of
the application of a rule, an history is built from state of an orbit at the
moment of a rule application (bottom of a DAG) up to its origins (top of a
DAG). %
Moreover, there are two kind of layers in an evaluation DAG: First layer is an
event layer, indexed by a rule's application number from the persistent name and
showing a list of topological changes provoked by the rule; Second layer is an
orbit level, indexed by a node (related to the application number) from the
persistent name and showing a list of the orbits produced by the topological
changes of event level above. %

For example, $PN_{3}$'s evaluation DAG (Fig.~\ref{fig:eval-dag-pn3}) references
an orbit which the extrude rule expects to be a face $\orbit{0,1}$. %
Considering $PN_3=[1n_{5};2n_{6};3n_{0}]$, its evaluation DAG starts with the
application $3$: The orbit level is indexed by $n_{0}$ and contains a face orbit
$\orbit{0,1}$; The event level above is indexed with the application number
$3$. %
A query to a bulletin board with an orbit $\orbit{0,1}$ incident to the node
$n_{0}$ will find that the orbit $\orbit{0,1}$ incident to $n_{0}$ was modified
by the application $3$ so the event listed in the level event is
$MODIFICATION$. %
Continuing with application $2$: The orbit level is indexed with $n_{4}$ and
since the face orbit was modified (\textit{i.e.} no creation nor split nor
merging of the orbit), it can be traced up to this level; The event level, this
time, is indexed by $2$ and a query will find that the face orbit was created by
this application. %
Finally, …


\subsection{Traces and origins}

\begin{figure}[ht!]
  \centering
  \includegraphics[width=0.4\textwidth]{pn7}
  \caption{Evaluation DAG of the Persistent Name $PN_{7}$}
  \label{fig:eval-dag-pn7}
\end{figure}

\subsection{Paths to origins}

\begin{figure}[ht!]
  \centering
  \includegraphics[width=0.4\textwidth]{pn6}
  \caption{Evaluation DAG of the Persistent Name $PN_{6}$}
  \label{fig:eval-dag-pn6}
\end{figure}

%%% Local Variables:
%%% mode: latex
%%% TeX-master: "../main"
%%% End:
